\chapter{Derivation of Perihelion Precession from Orbital Strain Coupling}
\label{app:perihelion_orbital_strain}
\addtocontents{toc}{\protect\setcounter{tocdepth}{0}}

\noindent\textbf{Subject:} Deriving the weak-field perihelion orbital correction from conservative orbital strain coupling in the active Soft Floor.

This appendix derives the perihelion orbital correction from Tension Field Theory (TFT) using a mechanical specific-strain-potential argument. The purpose is to show that the standard weak-field perihelion equation of General Relativity (GR) can be recovered without treating spacetime curvature as the fundamental cause.

In General Relativity, perihelion precession is described as the post-Newtonian effect of curved spacetime on a massive orbit. In Tension Field Theory, the same effect is interpreted as the accumulated orbital phase shift of a bound three-dimensional topological defect moving tangentially through the radial strain field of the active FCC Soft Floor.

This derivation does not duplicate the optical derivations of gravitational redshift, Shapiro delay, or gravitational lensing. Those effects concern wave propagation through a strain-induced refractive index. Perihelion precession is different. It is a bound-mass orbital effect. Therefore, it must be derived from the mechanical coupling between radial tension strain and tangential orbital transport.

\section{Newtonian Closed-Orbit Baseline}

A central mass $M$ produces the ordinary radial gravitational acceleration
\begin{equation}
    g_{\mathrm{b}} = \frac{GM}{r^2}.
\end{equation}

The corresponding inward radial acceleration is
\begin{equation}
    a_r = - \frac{GM}{r^2}.
\end{equation}

In Newtonian mechanics, this exact inverse-square acceleration produces a closed Keplerian ellipse. Define
\begin{equation}
    u = \frac{1}{r}.
\end{equation}

The Newtonian orbital equation is
\begin{equation}
    \frac{\mathrm{d}^2u}{\mathrm{d}\varphi^2} + u = \frac{GM}{\ell_{\mathrm{orb}}^2},
\end{equation}
where $\ell_{\mathrm{orb}}$ is the specific orbital angular momentum of the orbiting body:
\begin{equation}
    \ell_{\mathrm{orb}} = rv_\perp.
\end{equation}
Here $v_\perp$ is the tangential orbital velocity.

The Newtonian solution is
\begin{equation}
    u(\varphi) = \frac{1}{p} \left( 1+e\cos\varphi \right),
\end{equation}
where
\begin{equation}
    p = \frac{\ell_{\mathrm{orb}}^2}{GM} = a(1-e^2).
\end{equation}
Here $a$ is the semi-major axis and $e$ is the orbital eccentricity. Because the angular dependence is exactly $\cos\varphi$, the orbit closes after one full angular cycle,
\begin{equation}
    \varphi = 2\pi.
\end{equation}

Therefore, a pure inverse-square radial force produces no perihelion precession. Any perihelion shift must arise from an additional correction to the Newtonian radial law.

\section{Radial Soft Floor Strain Potential}

In TFT, the central mass does not curve an empty manifold at the fundamental level. Instead, it imposes a radial strain on the active Soft Floor.

The dimensionless gravitational strain is
\begin{equation}
    \alpha_{\mathrm{G}}(r) = \frac{GM}{rc_{\mathrm{E}}^2}.
\end{equation}

The corresponding radial strain potential per unit mass is
\begin{equation}
    \Phi_\tau(r) = - c_{\mathrm{E}}^2 \alpha_{\mathrm{G}}(r).
\end{equation}

Substituting the expression for $\alpha_{\mathrm{G}}$,
\begin{equation}
    \Phi_\tau(r) = - c_{\mathrm{E}}^2 \left( \frac{GM}{rc_{\mathrm{E}}^2} \right).
\end{equation}

Therefore,
\begin{equation}
    \Phi_\tau(r) = - \frac{GM}{r}.
\end{equation}

This is algebraically identical to the Newtonian potential, but its TFT interpretation is different. In TFT, $\Phi_\tau$ is the radial strain potential of the active Soft Floor, not the curvature of empty spacetime.

Taking its radial derivative gives the ordinary inward acceleration:
\begin{equation}
    a_\tau = - \frac{\mathrm{d}\Phi_\tau}{\mathrm{d}r}.
\end{equation}

Since
\begin{equation}
    \frac{\mathrm{d}}{\mathrm{d}r} \left( - \frac{GM}{r} \right) = \frac{GM}{r^2},
\end{equation}
we obtain
\begin{equation}
    a_\tau = - \frac{GM}{r^2}.
\end{equation}

The inward acceleration magnitude is therefore
\begin{equation}
    g_{\mathrm{b}} = \frac{GM}{r^2}.
\end{equation}

This recovers the Newtonian closed-orbit limit.

\section{Tangential Loading of a Bound Mass Defect}

A planet is not a static test mass inside the radial strain field. It is a bound topological-defect aggregate moving tangentially through that strain field.

Let the tangential orbital speed be
\begin{equation}
    v_\perp.
\end{equation}

The dimensionless tangential loading of the Soft Floor is
\begin{equation}
    q_{\mathrm{orb}} = \frac{v_\perp^2}{c_{\mathrm{E}}^2}.
\end{equation}

This quantity measures the fraction of the Soft Floor's kinematic capacity consumed by tangential orbital transport.

In TFT, this tangential loading is mechanically significant because a massive body is not a massless wave. A photon follows a refractive path through the strain field, but a massive body must preserve its three-dimensional bound topology while moving sideways through the radial tension gradient.

Therefore, the orbiting body produces a local conservative cross-coupling between radial strain and tangential transport.

\section{Mechanical Origin of the Orbital Coupling Term}

To understand why the specific orbital strain potential takes the form
\begin{equation}
    \Phi_{\mathrm{orb}}
    =
    \Phi_\tau
    \frac{v_\perp^2}{c_{\mathrm{E}}^2},
\end{equation}
we must treat the active vacuum as a pre-stressed nonlinear elastic lattice rather than as empty geometry.

A central mass establishes a radial strain potential in the Soft Floor,
\begin{equation}
    \Phi_\tau
    =
    -
    \frac{GM}{r}.
\end{equation}

This radial potential is not merely a passive background. It is a real pre-existing strain state of the active FCC lattice.

When a bound mass defect moves tangentially through this radial strain field, it must continuously preserve its three-dimensional topological closure while being transported orthogonally to the radial tension gradient. The Soft Floor must therefore perform additional conservative structural work to maintain the defect's bound configuration during orbital motion.

In continuum mechanics, a pre-stressed medium responds differently to an additional deformation than an unstressed medium. A deformation applied orthogonally to an existing prestress produces an incremental strain-energy density proportional to the product of the existing prestress and the newly applied strain:
\begin{equation}
    \text{incremental strain-energy density}
    \sim
    \text{prestress}
    \times
    \text{incremental strain}.
\end{equation}

After normalizing the corresponding strain-energy density by the inertial mass density of the bound defect, this prestress interaction appears as a specific strain potential per unit mass.

In the TFT orbital case, the radial prestress is represented by the specific strain potential
\begin{equation}
    \Phi_\tau.
\end{equation}

The incremental tangential strain is represented by the fraction of the Soft Floor kinematic capacity consumed by orbital motion:
\begin{equation}
    \frac{v_\perp^2}{c_{\mathrm{E}}^2}.
\end{equation}

This quantity is not the Newtonian kinetic energy per unit mass. It is a dimensionless structural loading factor. It measures how strongly tangential motion consumes the available elastic propagation capacity of the lattice. Therefore, no classical kinetic-energy factor of \(1/2\) appears in this term.

The leading conservative orbital strain potential is therefore
\begin{equation}
    \Phi_{\mathrm{orb}}
    =
    \Phi_\tau
    \left(
    \frac{v_\perp^2}{c_{\mathrm{E}}^2}
    \right).
\end{equation}

This is the central constitutive relation of the perihelion derivation. It states that transverse motion through a radially pre-strained Soft Floor generates an additional conservative strain-locking potential per unit mass.

Physically, the term means
\begin{equation}
    \text{radial Soft Floor strain}
    +
    \text{tangential transport of a bound defect}
    \rightarrow
    \text{orbital strain-locking potential}.
\end{equation}

This mechanism is distinct from dark-matter-like cosmic impedance. It is the local conservative correction required for a bound three-dimensional topological defect to maintain structural closure while orbiting through a radial lattice strain field.

\section{Cubic Reduced Orbital Strain Potential}

Substitute
\begin{equation}
    \Phi_\tau
    =
    -
    \frac{GM}{r}
\end{equation}
into the orbital coupling term:
\begin{equation}
    \Phi_{\mathrm{orb}}
    =
    -
    \frac{GM}{r}
    \frac{v_\perp^2}{c_{\mathrm{E}}^2}.
\end{equation}

The specific orbital angular momentum is conserved for the weak-field spherical bound-orbit problem:
\begin{equation}
    \ell_{\mathrm{orb}}
    =
    rv_\perp.
\end{equation}

\noindent\textit{Reduction note.}
At this step, the orbital strain term is being reduced to the central-force bound-orbit problem by imposing conservation of specific angular momentum. Therefore, the potential used below should be read as a reduced central-orbit strain potential, not as a fully general velocity-dependent field potential for arbitrary motion. The derivation applies to weak-field, spherical, conservative bound orbits.

Therefore,
\begin{equation}
    v_\perp
    =
    \frac{\ell_{\mathrm{orb}}}{r},
\end{equation}
and
\begin{equation}
    v_\perp^2
    =
    \frac{\ell_{\mathrm{orb}}^2}{r^2}.
\end{equation}

Substituting this into the orbital strain potential gives
\begin{equation}
    \Phi_{\mathrm{orb}}^{(\mathrm{red})}
    =
    -
    \frac{GM}{r}
    \frac{\ell_{\mathrm{orb}}^2}{r^2c_{\mathrm{E}}^2}.
\end{equation}

Therefore,
\begin{equation}
    \Phi_{\mathrm{orb}}^{(\mathrm{red})}
    =
    -
    \frac{GM\ell_{\mathrm{orb}}^2}
    {c_{\mathrm{E}}^2r^3}.
\end{equation}

This is the key reduced mechanical result. The correction potential per unit mass scales as
\begin{equation}
    \Phi_{\mathrm{orb}}^{(\mathrm{red})}
    \sim
    -
    \frac{1}{r^3}.
\end{equation}

The cubic radial dependence arises because the correction combines the radial strain potential,
\begin{equation}
    \Phi_\tau
    \sim
    -
    \frac{1}{r},
\end{equation}
with the tangential orbital loading,
\begin{equation}
    v_\perp^2
    \sim
    \frac{1}{r^2}.
\end{equation}

Thus,
\begin{equation}
    \frac{1}{r}
    \times
    \frac{1}{r^2}
    =
    \frac{1}{r^3}.
\end{equation}

In TFT language, perihelion precession arises from this cubic reduced orbital strain-potential coupling.

\section{Radial Derivative and Emergence of the Coefficient \(3\)}

The reduced central-orbit strain potential is
\begin{equation}
    \Phi_{\mathrm{strain}}^{(\mathrm{red})}
    =
    \Phi_\tau
    +
    \Phi_{\mathrm{orb}}^{(\mathrm{red})}.
\end{equation}

Therefore,
\begin{equation}
    \Phi_{\mathrm{strain}}^{(\mathrm{red})}
    =
    -
    \frac{GM}{r}
    -
    \frac{GM\ell_{\mathrm{orb}}^2}
    {c_{\mathrm{E}}^2r^3}.
\end{equation}

The inward central acceleration generated by this reduced strain potential is obtained from
\begin{equation}
    a_{\mathrm{strain}}^{(\mathrm{red})}
    =
    -
    \frac{\mathrm{d}\Phi_{\mathrm{strain}}^{(\mathrm{red})}}{\mathrm{d}r}.
\end{equation}

The Newtonian strain term gives
\begin{equation}
    -
    \frac{\mathrm{d}}{\mathrm{d}r}
    \left(
    -
    \frac{GM}{r}
    \right)
    =
    -
    \frac{GM}{r^2}.
\end{equation}

The reduced orbital strain term gives
\begin{equation}
    -
    \frac{\mathrm{d}}{\mathrm{d}r}
    \left(
    -
    \frac{GM\ell_{\mathrm{orb}}^2}
    {c_{\mathrm{E}}^2r^3}
    \right)
    =
    -
    \frac{3GM\ell_{\mathrm{orb}}^2}
    {c_{\mathrm{E}}^2r^4}.
\end{equation}

Therefore, the reduced inward central acceleration generated by the strain field is
\begin{equation}
    a_{\mathrm{strain}}^{(\mathrm{red})}
    =
    -
    \left[
    \frac{GM}{r^2}
    +
    \frac{3GM\ell_{\mathrm{orb}}^2}
    {c_{\mathrm{E}}^2r^4}
    \right].
\end{equation}

Equivalently, the inward central acceleration magnitude used in the reduced orbital equation is
\begin{equation}
    g_{\mathrm{strain}}^{(\mathrm{red})}(r)
    =
    \frac{GM}{r^2}
    +
    \frac{3GM\ell_{\mathrm{orb}}^2}
    {c_{\mathrm{E}}^2r^4}.
\end{equation}

Since
\begin{equation}
    g_{\mathrm{b}}
    =
    \frac{GM}{r^2},
\end{equation}
and
\begin{equation}
    v_\perp^2
    =
    \frac{\ell_{\mathrm{orb}}^2}{r^2},
\end{equation}
the correction may also be written as
\begin{equation}
    g_{\mathrm{orb}}
    =
    3g_{\mathrm{b}}
    \frac{v_\perp^2}{c_{\mathrm{E}}^2}.
\end{equation}

This expression should be understood as the reduced bound-orbit correction, not as an independent static gravitational field term.

The coefficient \(3\) is therefore not inserted as a numerical fitting factor. It emerges from the radial derivative of the cubic reduced strain-potential term:
\begin{equation}
    \frac{\mathrm{d}}{\mathrm{d}r}
    \left(
    r^{-3}
    \right)
    =
    -3r^{-4}.
\end{equation}

Thus, the origin of the coefficient is mechanical:
\begin{equation}
\begin{split}
    \text{radial strain potential} &\times \text{tangential orbital loading} \\
    &\rightarrow r^{-3} \text{ reduced strain potential} \\
    &\rightarrow 3r^{-4} \text{ bound-orbit acceleration correction}.
\end{split}
\end{equation}

\section{Conversion to the Orbital Equation}

For the reduced weak-field central-orbit problem, the inward acceleration magnitude used in Binet's equation is
\begin{equation}
    g_{\mathrm{strain}}^{(\mathrm{red})}(r)
    =
    \frac{GM}{r^2}
    +
    \frac{3GM\ell_{\mathrm{orb}}^2}
    {c_{\mathrm{E}}^2r^4}.
\end{equation}

For a reduced central-acceleration problem, the Binet equation is
\begin{equation}
    \frac{\mathrm{d}^2u}{\mathrm{d}\varphi^2}
    +
    u
    =
    \frac{
    g_{\mathrm{strain}}^{(\mathrm{red})}(1/u)
    }
    {\ell_{\mathrm{orb}}^2u^2},
\end{equation}
where
\begin{equation}
    u
    =
    \frac{1}{r}.
\end{equation}

Using
\begin{equation}
    \frac{1}{r^2} = u^2,
\end{equation}
and
\begin{equation}
    \frac{1}{r^4} = u^4,
\end{equation}
we obtain
\begin{equation}
    g_{\mathrm{strain}}^{(\mathrm{red})}(1/u) = GMu^2 + \frac{3GM\ell_{\mathrm{orb}}^2}{c_{\mathrm{E}}^2} u^4.
\end{equation}

Substituting this into the Binet equation gives
\begin{equation}
    \frac{\mathrm{d}^2u}{\mathrm{d}\varphi^2} + u = \frac{ GMu^2 + \frac{3GM\ell_{\mathrm{orb}}^2}{c_{\mathrm{E}}^2}u^4 }{\ell_{\mathrm{orb}}^2u^2}.
\end{equation}

Separate the two terms:
\begin{equation}
    \frac{\mathrm{d}^2u}{\mathrm{d}\varphi^2} + u = \frac{GM}{\ell_{\mathrm{orb}}^2} + \frac{3GM}{c_{\mathrm{E}}^2}u^2.
\end{equation}

Thus, TFT recovers the standard weak-field post-Newtonian orbital equation:
\begin{equation}
    \boxed{ \frac{\mathrm{d}^2u}{\mathrm{d}\varphi^2} + u = \frac{GM}{\ell_{\mathrm{orb}}^2} + \frac{3GM}{c_{\mathrm{E}}^2}u^2 }.
\end{equation}

This is the same orbital equation obtained in General Relativity for the weak-field perihelion correction.

\section{Recovery of Perihelion Precession}

The Newtonian solution is
\begin{equation}
    u_N(\varphi)
    =
    \frac{1}{p}
    \left(
    1+e\cos\varphi
    \right),
\end{equation}
where
\begin{equation}
    p
    =
    \frac{\ell_{\mathrm{orb}}^2}{GM}
    =
    a(1-e^2).
\end{equation}

The TFT correction adds the nonlinear term
\begin{equation}
    \frac{3GM}{c_{\mathrm{E}}^2}u^2.
\end{equation}

To first order, substituting the Newtonian solution
\begin{equation}
    u_N(\varphi)
    =
    \frac{1}{p}
    \left(
    1+e\cos\varphi
    \right)
\end{equation}
into the nonlinear correction produces a resonant term proportional to \(e\cos\varphi\). This resonant term is absorbed as a small shift of the orbital angular frequency.

The corrected orbit may therefore be written approximately as
\begin{equation}
    u(\varphi)
    \approx
    \frac{1}{p}
    \left[
    1+
    e\cos\left((1-\delta)\varphi\right)
    \right],
\end{equation}
where the small precession parameter is
\begin{equation}
    \delta
    =
    \frac{3GM}{pc_{\mathrm{E}}^2}.
\end{equation}

The perihelion advance per orbit is
\begin{equation}
    \Delta\varphi
    =
    2\pi\delta.
\end{equation}

Substituting \(\delta\) gives
\begin{equation}
    \Delta\varphi
    =
    2\pi
    \left(
    \frac{3GM}{pc_{\mathrm{E}}^2}
    \right).
\end{equation}

Therefore,
\begin{equation}
    \Delta\varphi
    =
    \frac{6\pi GM}
    {pc_{\mathrm{E}}^2}.
\end{equation}

Using
\begin{equation}
    p
    =
    a(1-e^2),
\end{equation}
we obtain
\begin{equation}
    \boxed{
    \Delta\varphi
    =
    \frac{6\pi GM}
    {a(1-e^2)c_{\mathrm{E}}^2}
    }.
\end{equation}

This is the standard weak-field perihelion precession formula. In TFT, the same orbital phase advance is interpreted as the cumulative effect of the cubic reduced Soft Floor strain potential rather than as fundamental spacetime curvature.

\section{Mechanical Interpretation}

In General Relativity, the non-closing orbit is attributed to the curvature of spacetime around the central mass. In TFT, the same mathematical correction is interpreted differently.

The central mass creates a radial strain potential in the active Soft Floor:
\begin{equation}
    \Phi_\tau
    =
    -
    \frac{GM}{r}.
\end{equation}

The orbiting mass is a bound three-dimensional topological defect that must preserve its structural closure while moving tangentially through this radial strain field.

This produces an additional conservative specific orbital strain-potential coupling:
\begin{equation}
    \Phi_{\mathrm{orb}}
    =
    \Phi_\tau
    \frac{v_\perp^2}{c_{\mathrm{E}}^2}.
\end{equation}

After imposing conservation of specific angular momentum,
\begin{equation}
    \ell_{\mathrm{orb}}
    =
    rv_\perp,
\end{equation}
this coupling reduces to the central-orbit strain potential
\begin{equation}
    \Phi_{\mathrm{orb}}^{(\mathrm{red})}
    =
    -
    \frac{GM\ell_{\mathrm{orb}}^2}
    {c_{\mathrm{E}}^2r^3}.
\end{equation}

Because the reduced tangential orbital loading scales as
\begin{equation}
    v_\perp^2
    \sim
    r^{-2},
\end{equation}
and the radial strain potential scales as
\begin{equation}
    \Phi_\tau
    \sim
    r^{-1},
\end{equation}
their reduced coupling scales as
\begin{equation}
    \Phi_{\mathrm{orb}}^{(\mathrm{red})}
    \sim
    r^{-3}.
\end{equation}

The radial derivative of this cubic reduced strain potential generates the additional inward bound-orbit correction:
\begin{equation}
    g_{\mathrm{orb}}
    =
    3g_{\mathrm{b}}
    \frac{v_\perp^2}{c_{\mathrm{E}}^2}.
\end{equation}

This expression is not an independent static gravitational field term. It is the reduced conservative correction experienced by a massive bound orbit after angular momentum conservation has been imposed.

This correction slightly over-closes the radial restoring force during each orbit, causing the ellipse to fail to return to its original orientation. The result is perihelion precession.

Thus,
\begin{equation}
    \boxed{
    \text{perihelion precession}
    =
    \text{cumulative orbital phase advance from cubic reduced Soft Floor strain potential}
    }.
\end{equation}

\section{Distinction from the Dark-Matter Impedance Term}

The perihelion correction derived in this appendix must be distinguished from the galactic impedance mechanism developed in \hyperref[app:cosmic_impedance]{Appendix~\ref*{app:cosmic_impedance}}.

The dark-matter-like acceleration correction is associated with large-scale inductive coupling of the interstitial condensate:

\begin{equation}
    g_{\mathrm{DM}}
    =
    \Xi(R)\sqrt{a_{\mathrm{D}}g_{\mathrm{b}}}.
\end{equation}

This term is an environmental response of the active vacuum. It becomes relevant in the low-acceleration galactic regime,

\begin{equation}
    g_{\mathrm{b}}
    \ll
    a_{\mathrm{D}},
\end{equation}

where the interstitial condensate contributes an additional large-scale impedance response to baryonic gravity.

Solar-system perihelion precession occurs in the opposite regime,

\begin{equation}
    g_{\mathrm{b}}
    \gg
    a_{\mathrm{D}}.
\end{equation}

In this high-acceleration regime, the galactic impedance term is screened or negligible:

\begin{equation}
    g_{\mathrm{DM}}
    \rightarrow
    0.
\end{equation}

The perihelion correction is therefore not a dark-matter effect. It is a local conservative bound-orbit correction produced by tangential transport of a massive defect through a radial Soft Floor strain field.

The environmental acceleration hierarchy may be written schematically as

\begin{equation}
    g_{\mathrm{env}}
    =
    g_{\mathrm{b}}
    +
    g_{\mathrm{DM}},
\end{equation}

where

\begin{equation}
    g_{\mathrm{DM}}
    =
    \Xi(R)\sqrt{a_{\mathrm{D}}g_{\mathrm{b}}}.
\end{equation}

For precision massive bound-orbit dynamics, this environmental field is supplemented by the reduced orbital correction

\begin{equation}
    g_{\mathrm{bound}}
    =
    g_{\mathrm{env}}
    +
    g_{\mathrm{orb}},
\end{equation}

where

\begin{equation}
    g_{\mathrm{orb}}
    =
    3g_{\mathrm{b}}
    \frac{v_\perp^2}{c_{\mathrm{E}}^2}.
\end{equation}

Thus, \(g_{\mathrm{DM}}\) and \(g_{\mathrm{orb}}\) arise from different mechanical channels. The term \(g_{\mathrm{DM}}\) is a large-scale condensate impedance response, while \(g_{\mathrm{orb}}\) is a reduced conservative correction for massive bound orbits. In the solar system, \(g_{\mathrm{DM}}\) is negligible while \(g_{\mathrm{orb}}\) produces the weak-field perihelion correction. In galactic outskirts, \(g_{\mathrm{orb}}\) is negligible while \(g_{\mathrm{DM}}\) dominates the non-Newtonian rotation-curve behavior.

\section{Relation to the Effective GR Metric}

The same result may be expressed in effective metric language. The perihelion correction recovered here is the standard first post-Newtonian orbital correction associated with the weak-field Schwarzschild limit.

In the parametrized post-Newtonian description, the standard perihelion advance is governed by the combination of metric coefficients that reduces to the General Relativistic value when
\begin{equation}
    \beta_{\mathrm{PPN}} = 1,
    \qquad
    \gamma_{\mathrm{PPN}} = 1.
\end{equation}

The present appendix does not independently derive the full PPN metric. Instead, it shows that the TFT reduced bound-orbit strain potential generates the same weak-field orbital equation and therefore the same perihelion advance as the Schwarzschild weak-field limit.

The correspondence is therefore
\begin{equation}
    \text{GR weak-field orbital curvature correction}
    \quad
    \longleftrightarrow
    \quad
    \text{TFT reduced conservative orbital strain correction}.
\end{equation}

The distinction is ontological: GR represents the correction as spacetime curvature, while TFT represents it as conservative strain-locking of a bound topological defect moving through the radially pre-stressed Soft Floor.

\section{Conclusion}

TFT reproduces the GR perihelion orbital equation from a conservative reduced orbital strain-potential correction.

Starting from the radial Soft Floor strain potential,
\begin{equation}
    \Phi_\tau = - \frac{GM}{r},
\end{equation}
and coupling it to tangential orbital loading,
\begin{equation}
    \frac{v_\perp^2}{c_{\mathrm{E}}^2},
\end{equation}
yields the cubic reduced specific orbital strain potential
\begin{equation}
    \Phi_{\mathrm{orb}}^{(\mathrm{red})}
=
-
\frac{GM\ell_{\mathrm{orb}}^2}
{c_{\mathrm{E}}^2r^3}.
\end{equation}

The radial derivative of this term generates the acceleration correction
\begin{equation}
    g_{\mathrm{orb}} = 3g_{\mathrm{b}} \frac{v_\perp^2}{c_{\mathrm{E}}^2}.
\end{equation}

This leads directly to the orbital equation
\begin{equation}
    \frac{\mathrm{d}^2u}{\mathrm{d}\varphi^2} + u = \frac{GM}{\ell_{\mathrm{orb}}^2} + \frac{3GM}{c_{\mathrm{E}}^2}u^2,
\end{equation}
and therefore to the perihelion advance
\begin{equation}
    \Delta\varphi = \frac{6\pi GM}{a(1-e^2)c_{\mathrm{E}}^2}.
\end{equation}

Thus, the standard GR perihelion result is recovered, while its interpretation is shifted from fundamental spacetime curvature to the conservative strain-locking dynamics of a bound topological defect moving through the active vacuum lattice.